\documentclass[11pt,a4paper]{article}
\usepackage[T1]{fontenc}
\usepackage[utf8]{inputenc}
\usepackage[english]{babel}
\usepackage{amsmath,amssymb}
\usepackage[margin=2cm]{geometry}
\setlength{\parindent}{0pt}
\setlength{\emergencystretch}{3em}
\allowdisplaybreaks

\title{About \texttt{abstract\_structures.tex}}
\author{}
\date{}
\begin{document}
\maketitle
This note explains what the test data in \texttt{abstract\_structures.tex} is used for. That file is input for the engine in \texttt{main.c}; it is read by the test \texttt{test\_algebraic\_structures.c} in this folder.

\begin{enumerate}
\item \textbf{Purpose.} The file holds 24 true statements about an arbitrary set with an operation, maps, a subset or a relation of which only the hypotheses are known. Each must be proved, with a written proof.

\item \textbf{What it exercises.} The language of paragraph 17o of \texttt{parser.c} and chapter 28 of \texttt{main.c}: proof by contradiction with instances of the hypotheses, substitution of equals (congruence closure), named witnesses, cases; every written step is re-checked, and no model with few elements may contradict the statement. The two statements that begin with $\exists S$ are proved by exhibiting a finite model.

\item \textbf{Variety.} Uniqueness of the neutral element, left and right inverses, generalised associativity, cancellation in a group, an operation with $x \perp x = e$ is commutative; compositions of injective and of surjective maps, a left inverse, an involution; pairs; relations that are symmetric, transitive, euclidean, with quantifiers written without a set; subsets closed under maps; a statement that needs cases.

\item \textbf{Expected result.} The line \texttt{PASS: 24 of 24 statements proved, each with a written proof}.

\item \textbf{How to run it.} From the top folder of the repository:
\begin{verbatim}
cc -std=c11 -O2 -Wall -Wextra -o /tmp/test_algebraic_structures \
    tests/algebraic_structures/test_algebraic_structures.c \
    && /tmp/test_algebraic_structures
\end{verbatim}
\end{enumerate}
\end{document}
